% Standalone typesetting reference. All numbers and graphics are synthetic.
% This file deliberately uses standard packages and also works with pdfLaTeX.
\documentclass[11pt,a4paper]{article}
\usepackage[margin=25mm]{geometry}
\usepackage{amsmath,amssymb,graphicx,booktabs,array,tabularx,longtable}
\usepackage{siunitx,caption,subcaption,pdflscape}
\usepackage[numbers,sort&compress]{natbib}
\usepackage[hidelinks]{hyperref}
\newcolumntype{Y}{>{\raggedright\arraybackslash}X}
\sisetup{uncertainty-mode=separate,group-minimum-digits=5}
\captionsetup{font=small,labelfont=bf}
\setlength{\emergencystretch}{2em}
\setlength{\parskip}{5pt}
\setlength{\parindent}{0pt}
\title{Prism figure, table and equation examples}
\author{Typesetting examples only}
\date{}

\begin{document}
\maketitle
These examples demonstrate layout, cross-references and pagination. Every
number, label and drawing is synthetic and carries no scientific claim.
Copy only the relevant environment into a manuscript. Compile this reference
independently with \texttt{--main examples.tex --no-build-script}.

\section{Related figure panels}
Figure~\ref{fig:grid} demonstrates a four-panel arrangement. Each panel is a
separate reference target; for example, Fig.~\ref{fig:grid-b}. In a real paper,
keep panel lettering and graphic scales consistent and explain the uncertainty
in the main legend.

\begin{figure}[htbp]
\centering
\begin{subfigure}[t]{0.47\linewidth}
\centering
\setlength{\unitlength}{1pt}
\begin{picture}(165,100)
\put(20,18){\vector(1,0){135}}\put(20,18){\vector(0,1){72}}
\put(37,18){\line(0,1){20}}\put(64,18){\line(0,1){42}}
\put(91,18){\line(0,1){48}}\put(118,18){\line(0,1){61}}
\put(60,0){\makebox(60,12){\small Example setting}}
\put(25,86){\makebox(110,12){\small Synthetic comparison}}
\end{picture}
\caption{A bar-like comparison.}\label{fig:grid-a}
\end{subfigure}\hfill
\begin{subfigure}[t]{0.47\linewidth}
\centering
\setlength{\unitlength}{1pt}
\begin{picture}(165,100)
\put(20,18){\vector(1,0){135}}\put(20,18){\vector(0,1){72}}
\qbezier(25,25)(75,85)(145,76)
\put(35,38){\circle*{4}}\put(65,60){\circle*{4}}
\put(95,74){\circle*{4}}\put(135,77){\circle*{4}}
\put(60,0){\makebox(60,12){\small Example setting}}
\put(25,86){\makebox(110,12){\small Synthetic trend}}
\end{picture}
\caption{A trend with markers.}\label{fig:grid-b}
\end{subfigure}

\medskip
\begin{subfigure}[t]{0.47\linewidth}
\centering
\fbox{\parbox[c][27mm][c]{0.91\linewidth}{\centering
Input $\longrightarrow$ transformation\\[5pt]
$\longrightarrow$ output}}
\caption{A simple process.}\label{fig:grid-c}
\end{subfigure}\hfill
\begin{subfigure}[t]{0.47\linewidth}
\centering
\fbox{\parbox[c][27mm][c]{0.91\linewidth}{\centering
Replace this box with a PDF or PNG\\[5pt]
\texttt{\string\includegraphics}\\[3pt]
\texttt{[width=\string\linewidth]\{file\}}}}
\caption{An external image slot.}\label{fig:grid-d}
\end{subfigure}
\caption{\textbf{A related set of panels.} (a) Bars demonstrate equal plotting
areas. (b) Markers and a curve demonstrate distinct visual cues. (c) A diagram
demonstrates a process panel. (d) An image slot demonstrates a panel with a
different content type. The drawings are made entirely in TeX; replace them
with the study's own editable graphics. No data were analysed.}
\label{fig:grid}
\end{figure}
\clearpage

\section{A tall figure and a long legend}
Reserve space for both the graphic and its caption. Figure~\ref{fig:tall} uses
a graphic area of half the text height. Replace the slot with an image
bounded by \texttt{width=\string\linewidth},
\texttt{height=0.50\string\textheight} and \texttt{keepaspectratio}.
The bound should be checked against the actual legend and the final labels.

\begin{figure}[htbp]
\centering
\fbox{\parbox[c][0.50\textheight][c]{0.85\linewidth}{\centering
\Large Tall composite figure\\[12pt]
\normalsize Panel a: overview\\[26pt]
Panel b: detailed view\\[26pt]
Panel c: validation view\\[12pt]
\small Synthetic layout slots only}}
\caption{\textbf{A tall figure with a substantial legend.} The upper panel is
reserved for an overview that introduces the structure of the figure. The
middle panel is reserved for the detail needed to understand the main
comparison. The lower panel is reserved for a related validation view. In a
real manuscript, define all symbols and units, state what the markers and
lines mean, and describe the sample size and uncertainty supplied by the
study. Keep essential information in the caption so that the graphic can be
understood with its legend. Explain methodological choices in the Methods
section rather than repeating them here. This paragraph is deliberately long
enough to demonstrate how the legend consumes page space below a tall image.
It reports no experiment and must be replaced or removed before submission.}
\label{fig:tall}
\end{figure}
\clearpage

\section{Numerical values and uncertainty}
Table~\ref{tab:uncertainty} aligns values and their uncertainties without
scaling the table as an image. The missing value bypasses numerical parsing
with \texttt{\string\multicolumn}; define that convention in the table note.

\begin{table}[htbp]
\centering
\caption{Synthetic numerical examples, with no associated experiment.}
\label{tab:uncertainty}
\small
\begin{tabularx}{\linewidth}{@{}Y S[table-format=3.0]
  S[table-format=1.3(3)] S[table-format=-1.3]@{}}
\toprule
Setting & {Count} & {Estimate $\pm$ SD} & {Change}\\
\midrule
Reference setting & 120 & 0.385 +- 0.006 & -0.015\\
Alternative setting & 98 & 0.608 +- 0.051 & 0.208\\
Sensitivity setting & 36 & 0.517 +- 0.012 & -0.004\\
Unavailable setting & 0 & \multicolumn{1}{c}{---} & 0.000\\
\bottomrule
\end{tabularx}
\par\smallskip\footnotesize\raggedright
Notes. SD denotes the illustrative standard deviation. Counts, estimates and
changes are invented for layout only. The dash denotes an unavailable value,
not zero. Preserve the precision of the real source data when replacing cells.
\end{table}

\section{Equations that stay readable}
Use \texttt{align} for related equalities. Equation~\eqref{eq:example-gradient}
shows a reference to a specific line, and equation~\eqref{eq:example-split}
has one number for a multi-line expression.
\begin{align}
  \mathcal{L}(\theta) &= \frac{1}{n}\sum_{i=1}^{n}
    \ell\bigl(f_\theta(x_i),y_i\bigr)+\lambda\lVert\theta\rVert_2^2,
    \label{eq:example-loss}\\
  \nabla_\theta\mathcal{L}(\theta) &= \frac{1}{n}\sum_{i=1}^{n}
    \nabla_\theta\ell\bigl(f_\theta(x_i),y_i\bigr)+2\lambda\theta.
    \label{eq:example-gradient}
\end{align}
\begin{equation}
\begin{split}
  \Delta_t &= \mathcal{L}(\theta_{t+1})-\mathcal{L}(\theta_t)\\
  &\leq \langle\nabla\mathcal{L}(\theta_t),\theta_{t+1}-\theta_t\rangle
    +\frac{L}{2}\lVert\theta_{t+1}-\theta_t\rVert_2^2.
\end{split}\label{eq:example-split}
\end{equation}
These expressions demonstrate notation only. They are not a proof of a
property of a study. Define assumptions before adopting any mathematical claim.
\begin{subequations}\label{eq:example-system}
\begin{align}
 h(z)&=\begin{cases}z,&z\geq0,\\\alpha z,&z<0,\end{cases}
   \label{eq:example-cases}\\
 \mathbf{z}_{t+1}&=\begin{bmatrix}1&\eta\\0&1\end{bmatrix}\mathbf{z}_t
   +\begin{bmatrix}0\\u_t\end{bmatrix}.
   \label{eq:example-matrix}
\end{align}
\end{subequations}
The complete group is equation~\eqref{eq:example-system}; its two parts are
equations~\eqref{eq:example-cases} and~\eqref{eq:example-matrix}.
\clearpage

\section{A table that genuinely spans pages}
Table~\ref{tab:long-example} repeats its column headings and continuation
notice on each page. Do not wrap a \texttt{longtable} in a table float. Rows
below are invented layout checks rather than reported validation results.
\small
\begin{longtable}{@{}p{0.12\linewidth}p{0.25\linewidth}p{0.54\linewidth}@{}}
\caption{Synthetic inventory used to exercise page continuation.}
\label{tab:long-example}\\
\toprule Item & Setting & Explanation\\\midrule
\endfirsthead
\multicolumn{3}{@{}l}{Table~\thetable\ (continued)}\\[3pt]
\toprule Item & Setting & Explanation\\\midrule
\endhead
\midrule\multicolumn{3}{r@{}}{Continued on next page}\\
\endfoot
\bottomrule
\endlastfoot
C01 & Input identifier & Record a stable identifier and describe where the input originated.\\
C02 & Input version & Record the exact version rather than relying on a mutable label.\\
C03 & Units & Define the units for each numerical measurement in the source.\\
C04 & Missing values & Document which entries are unavailable and how they are represented.\\
C05 & Inclusion rule & State the rule used to select the material before analysis.\\
C06 & Exclusion rule & State the rule used to exclude material before analysis.\\
C07 & Partition & Describe the separation of training and evaluation material.\\
C08 & Identifier mapping & Preserve the mapping between source records and derived results.\\
C09 & Transformation & Describe any preprocessing that changes the scale of a variable.\\
C10 & Normalisation & State the reference quantities used for any normalisation.\\
C11 & Parameters & Record numerical settings together with the reason for their choice.\\
C12 & Initialisation & Record the initial condition or random seed where appropriate.\\
C13 & Repetition & Describe the unit of repetition without treating dependent samples as independent.\\
C14 & Sample count & State the exact number of observations relevant to each comparison.\\
C15 & Uncertainty & Define the uncertainty measure and the way it was calculated.\\
C16 & Interval & Identify whether intervals describe variability or estimation uncertainty.\\
C17 & Comparison & Define the reference condition for every reported difference.\\
C18 & Control & Describe the role of the control in the planned analysis.\\
C19 & Sensitivity & Describe how sensitive settings would be varied in a real study.\\
C20 & Validation & Describe an independent validation design where one is available.\\
C21 & Software & Identify the software and relevant version for each analysis stage.\\
C22 & Environment & Describe the computing environment when it affects reproducibility.\\
C23 & Hardware & Record hardware details when numerical performance depends on them.\\
C24 & Precision & Preserve the precision justified by the original measurements.\\
C25 & Rounding & Apply a documented display rounding rule without changing source values.\\
C26 & Aggregation & Describe the quantities being combined and the weighting used.\\
C27 & Dependence & Identify measurements sharing the same subject or source.\\
C28 & Error bars & Explain whether bars denote standard deviation, error or an interval.\\
C29 & Figure labels & Check that labels remain readable at their intended final dimensions.\\
C30 & Figure provenance & Retain the plotting source and the data used to create each figure.\\
C31 & Table provenance & Link table cells to the corresponding verified source results.\\
C32 & Equation symbols & Define variables, operators and dimensions before using them.\\
C33 & Assumptions & State the conditions under which a mathematical statement is used.\\
C34 & Cross-references & Check that figure, table and equation references resolve correctly.\\
C35 & Availability & Describe an access route and any actual restrictions accurately.\\
C36 & Limitations & Record limitations supported by the design and evidence of the study.\\
\end{longtable}
\normalsize

\begin{landscape}
\section{A wide table on a landscape page}
Use a dedicated page when many columns remain necessary after simplifying
headers. Table~\ref{tab:wide-example} keeps text upright on the rotated PDF
page and uses consistent cell sizes rather than shrinking the whole table.
\begin{table}[htbp]
\centering\small
\caption{A synthetic comparison with several explicitly labelled settings.}
\label{tab:wide-example}
\begin{tabularx}{\linewidth}{@{}Y *{8}{S[table-format=1.2]}@{}}
\toprule
Setting & {Metric A} & {Metric B} & {Metric C} & {Metric D}
  & {Metric E} & {Metric F} & {Metric G} & {Metric H}\\
\midrule
Reference configuration & 0.12 & 0.23 & 0.34 & 0.45 & 0.56 & 0.67 & 0.78 & 0.89\\
Alternative configuration & 0.21 & 0.32 & 0.43 & 0.54 & 0.65 & 0.76 & 0.87 & 0.98\\
Sensitivity configuration & 0.11 & 0.22 & 0.33 & 0.44 & 0.55 & 0.66 & 0.77 & 0.88\\
\bottomrule
\end{tabularx}
\par\smallskip\footnotesize\raggedright
All values are invented. Define each real metric, its unit and its direction
of interpretation in the accompanying manuscript.
\end{table}
\end{landscape}

\section{Bibliography mechanics}
The citation here demonstrates bibliography resolution only
\citep{wilkinson2016fair}; it is not evidence for any synthetic number above.
\bibliographystyle{unsrtnat}
\bibliography{references}
\end{document}
